\chapter{Reducing Stalls Through Dynamic Scheduling}
\label{ch:dynamic_scheduling}

In simple pipelines, instructions are fetched and issued strictly in program order. Unless a data dependency can be completely hidden by data forwarding, its presence will stall the entire pipeline until the dependency is cleared. Because instructions are executed in order, a stall prevents all subsequent instructions from proceeding, even if those later operations are entirely independent of the stalled data.

\section{The Limitation of In-Order Execution}
Consider the following sequence of floating-point instructions:
\begin{lstlisting}[language={[x86masm]Assembler}]
DIV.D   F0, F2, F4
ADD.D   F10, F0, F8
SUB.D   F12, F8, F14
\end{lstlisting}

In this example, the \texttt{ADD.D} instruction has a \gls{raw} hazard with the \texttt{DIV.D} instruction on register \texttt{F0}. In a standard in-order pipeline, the \texttt{ADD.D} must stall in the decode/issue stage until \texttt{DIV.D} completes its execution. Consequently, the \texttt{SUB.D} instruction is also blocked from executing, despite the fact that it only depends on \texttt{F8} and \texttt{F14}-registers that are completely independent of the ongoing division.

To resolve this inefficiency, processors employ \textbf{dynamic scheduling}, a technique that allows the hardware to rearrange the execution of instructions to minimize stalling.

\begin{definition}{Dynamic Scheduling}{dynamic-scheduling}
    Dynamic scheduling is a hardware-based technique in which the processor rearranges the execution of instructions \textit{at runtime} to reduce pipeline stalls, while strictly preserving the program's original \textbf{data flow} and \textbf{exception behavior}.
\end{definition}

This approach fundamentally implies \textbf{out-of-order execution} (and often out-of-order completion). While most commonly associated with complex superscalar architectures, dynamic scheduling can also be highly effective in simpler scalar processors.

\section{Advantages and Disadvantages}

Shifting the responsibility of instruction scheduling from the compiler (software) to the processor (hardware) introduces several structural tradeoffs.

\subsubsection*{Advantages}
\begin{itemize}
    \item \textbf{Runtime Dependency Resolution:} It can detect and handle dependencies that are unknown at compile time, such as dependencies dynamically generated by branch instructions or memory references resolving to the same address (aliasing).
    \item \textbf{Compiler Simplification:} It reduces the burden on the compiler to perfectly optimize instruction scheduling for every specific hardware iteration.
    \item \textbf{Binary Compatibility:} It decouples software from hardware, allowing the same compiled binary code to run efficiently across different processor architectures and pipeline depths without requiring recompilation.
\end{itemize}

\subsubsection*{Disadvantages}
\begin{itemize}
    \item \textbf{New Hazard Types:} Out-of-order execution reintroduces \gls{war} and \gls{waw} hazards, which are naturally avoided in strict in-order pipelines.
    \item \textbf{Hardware Complexity:} Implementing dynamic scheduling requires considerable additional hardware logic to track register states, buffered instructions, and dependency chains.
    \item \textbf{Imprecise Exceptions:} Exception handling becomes significantly more complex. A dynamically scheduled processor might generate an imprecise exception (where the processor state does not perfectly match the strict program order) because:
    \begin{itemize}
        \item Instructions that are \textit{later} in the program order than the faulting instruction might have already completed.
        \item Instructions that are \textit{earlier} in the program order than the faulting instruction might not have finished yet.
    \end{itemize}
\end{itemize}

\section{Implementation: Scoreboarding}

To implement dynamic scheduling, starting from a simple scalar architecture, we must provide a mechanism for instructions to execute out of order. This overtaking is facilitated by introducing multiple independent functional units rather than relying on a single, monolithic \texttt{EX} stage. Without multiple functional units, dynamic scheduling would be severely bottlenecked by structural hazards. 

Even with multiple functional units, instructions would normally still stall in the Instruction Decode (\texttt{ID}) phase due to data dependencies. To circumvent this, the hardware splits the traditional \texttt{ID} stage into two distinct phases: 
\begin{enumerate}
    \item \textbf{Issue:} Checks for structural hazards and issues the instruction to a functional unit.
    \item \textbf{Read Operands:} Waits for residual data dependencies to clear, and then reads the operands from the register file.
\end{enumerate}

This entire dependency-tracking and execution-control process is known as \textbf{Scoreboarding}.

Scoreboarding is a fundamental centralized approach to dynamic scheduling (first introduced in the CDC 6600). Instructions execute out of order only when there are sufficient hardware resources and no active data dependencies prevent it. The core of this mechanism is the \textbf{Scoreboard}, a hardware table that holds several instructions at a given time, dynamically checks for dependencies, and controls the execution of the functional units. Crucially, instructions are \textit{issued in strictly program order}, but they can bypass each other during the read operands phase, leading to out-of-order execution and completion.

\subsection{The Scoreboard Composition}

To accurately control the out-of-order execution pipeline, the Scoreboard must maintain comprehensive information about the state of the processor. The hardware tracking is divided into three primary tables:

\begin{itemize}
    \item \textbf{Instruction Status Table:} Tracks the current execution phase of each active instruction (Issue, Read Operands, Execute, or Write Result).
    \item \textbf{Functional Unit (FU) Status Table:} Tracks the state of all functional units in the system. Each entry contains several fields:
    \begin{itemize}
        \item \textbf{Busy:} A boolean flag indicating whether the unit is currently in use.
        \item \textbf{Op:} The specific operation to be performed in the unit (e.g., ADD, SUB, DIV).
        \item \textbf{Fi:} The destination register for the result.
        \item \textbf{Fj, Fk:} The source registers for the operation.
        \item \textbf{Qj, Qk:} The functional units that will produce the source registers \texttt{Fj} and \texttt{Fk} (if they are not yet ready).
        \item \textbf{Rj, Rk:} Boolean flags indicating whether the operands \texttt{Fj} and \texttt{Fk} are ready to be read.
    \end{itemize}
    \item \textbf{Register Result Status Table:} Indicates which functional unit will write a given register. If no active instruction is slated to write to a specific register, this field is left blank.
\end{itemize}

\subsection{Execution Phases}

Under Scoreboarding, an instruction's lifecycle through the processor pipeline is divided into four distinct phases, replacing the standard 5-stage RISC pipeline execution flow.

\subsubsection{1. Issue}
During this phase, the hardware decodes the instruction and checks for two specific hazards:
\begin{itemize}
    \item \textbf{Structural Hazard Check:} Is the required functional unit free (i.e., is its \texttt{Busy} flag set to \textit{false})?
    \item \textbf{\gls{waw} Hazard Check:} Is any other currently active instruction waiting to write to the same destination register (\texttt{Fi})? 
\end{itemize}
If both the required functional unit is available and no \gls{waw} hazards exist, the Scoreboard issues the instruction to the functional unit and updates its internal tables. If either a structural or a \gls{waw} hazard is present, the instruction \textit{stalls at Issue}, effectively blocking any subsequent instructions from issuing (since issue is strictly in-order).

\subsubsection{2. Read Operands}
Once issued, the instruction waits at the functional unit until its data operands become available, acting as the mechanism to resolve \gls{raw} hazards dynamically.
\begin{itemize}
    \item \textbf{\gls{raw} Hazard Check:} The Scoreboard monitors the \texttt{Rj} and \texttt{Rk} flags. If another active functional unit is currently computing a required operand, the flags remain false.
\end{itemize}
When both \texttt{Rj} and \texttt{Rk} evaluate to true, the required data dependencies are cleared. The Scoreboard commands the functional unit to read the operands from the register file and proceed. By waiting locally at the functional unit rather than in the centralized issue queue, independent instructions can bypass the waiting instruction.

\subsubsection{3. Execute}
Once the operands are read, the functional unit begins execution. The duration of this phase is dictated by the specific operation's latency (e.g., a floating-point multiplication will take more clock cycles than an integer addition). When the functional unit completes the computation, it explicitly notifies the Scoreboard that it is ready to write the result.

\subsubsection{4. Write Result}
Before the functional unit can write its newly computed value into the destination register (\texttt{Fi}), the Scoreboard must ensure that doing so will not overwrite a value that a previous instruction still needs to read. 
\begin{itemize}
    \item \textbf{\gls{war} Hazard Check:} The Scoreboard checks if any other active functional unit has this register listed as a source (\texttt{Fj} or \texttt{Fk}) but has \textit{not yet read it} (\texttt{Rj} or \texttt{Rk} is true, meaning they are still waiting to read the old value).
\end{itemize}
If a \gls{war} hazard is detected, the writing functional unit is stalled. It holds the computed result until all older instructions have successfully read the previous register value. Once the hazard is cleared, the result is written to the register file, the functional unit is freed, and the instruction execution formally completes.

\subsection{Example: Scoreboard Execution}

To illustrate how the Scoreboard manages out-of-order execution and prevents data corruption, consider the following sequence of instructions executing on a processor with multiple functional units:

\begin{lstlisting}[language={[x86masm]Assembler}]
LD      F6, 34(R2)
LD      F2, 45(R3)
MULTD   F0, F2, F4
SUBD    F8, F6, F2
DIVD    F10, F0, F6
ADDD    F6, F8, F2
\end{lstlisting}

Assume a classical CDC 6600 hardware configuration with single functional units per type and the following execution latencies:
\begin{itemize}
    \item \textbf{Load Unit (1 unit):} 1 clock cycle
    \item \textbf{FP Adder Unit (1 unit):} 2 clock cycles
    \item \textbf{FP Multiplier Unit (1 unit):} 10 clock cycles
    \item \textbf{FP Divider Unit (1 unit):} 40 clock cycles
\end{itemize}

\begin{table}[htpb]
    \centering
    \renewcommand{\arraystretch}{1.2}
    \begin{tabular}{l c c c c}
        \textbf{Instruction} & \textbf{Issue} & \textbf{Read Operands} & \textbf{Exec. Complete} & \textbf{Write Result} \\
        \hline
        \texttt{LD   F6, 34(R2)}  & 1  & 2  & 3  & 4  \\
        \texttt{LD   F2, 45(R3)}  & 5  & 6  & 7  & 8  \\
        \texttt{MULTD F0, F2, F4} & 6  & 9  & 19 & 20 \\
        \texttt{SUBD  F8, F6, F2} & 7  & 9  & 11 & 12 \\
        \texttt{DIVD  F10, F0, F6}& 8  & 21 & 61 & 62 \\
        \texttt{ADDD  F6, F8, F2} & 13 & 14 & 16 & 22 \\
        \hline
    \end{tabular}
    \caption{Instruction execution timing under single-unit functional constraints (CDC 6600 model).}
    \label{tab:scoreboard-timing}
\end{table}

As shown in \cref{tab:scoreboard-timing}, the Scoreboard dynamically orchestrates execution across the four pipeline phases:

\begin{itemize}
    \item \textbf{Structural Hazard Stalls at Issue:} The second \texttt{LD} cannot issue at cycle 2 because the single Load unit is occupied by the first \texttt{LD} until cycle 4; thus, it issues at cycle 5. Similarly, \texttt{ADDD} cannot issue at cycle 8 because the single FP Adder unit is occupied by \texttt{SUBD} until cycle 12, forcing \texttt{ADDD} to stall at the \textit{Issue} phase until cycle 13.
    \item \textbf{Out-of-Order Execution (\gls{raw} Handling):} \texttt{MULTD} issues at cycle 6 but stalls in \textit{Read Operands} until cycle 9, when \texttt{LD F2} finishes writing \texttt{F2} at cycle 8. Likewise, \texttt{DIVD} issues at cycle 8 but stalls in \textit{Read Operands} until cycle 21, waiting for \texttt{MULTD} to write \texttt{F0} at cycle 20. Meanwhile, \texttt{SUBD} reads its operands at cycle 9, completes execution at cycle 11, and writes back at cycle 12 bypassing both \texttt{MULTD} and \texttt{DIVD} completely.
    \item \textbf{Stalling on a \gls{war} Hazard:} The final \texttt{ADDD} instruction creates a critical \gls{war} hazard on register \texttt{F6} with \texttt{DIVD}. \texttt{DIVD} needs the old value of \texttt{F6} loaded by the first \texttt{LD}. Because \texttt{DIVD} is stalled waiting for \texttt{F0}, it only reads \texttt{F6} at cycle 21. \texttt{ADDD} finishes execution at cycle 16, but the Scoreboard strictly holds it in the \textit{Write Result} phase until cycle 22 immediately after \texttt{DIVD} performs its read-preventing register corruption.
\end{itemize}

\subsection{Limitations of Scoreboarding}

Although Scoreboarding successfully enables out-of-order execution, it suffers from several structural and microarchitectural bottlenecks that limit its overall effectiveness in modern high-performance processors:

\begin{itemize}
    \item \textbf{Lack of Data Forwarding:} Operands are read from the register file only when \textit{both} required source operands become available. The basic Scoreboard mechanism does not incorporate data forwarding (bypassing) paths, forcing dependent instructions to wait until results are formally written back to the register file before they can be read.
    \item \textbf{Non-Overlapping Pipeline Stages:} An additional clock cycle of latency is incurred because the \textit{Write Result} and \textit{Read Operands} stages cannot overlap. A register cannot be read in the exact same cycle that a computed result is being written to it.
    \item \textbf{Stalls on False Dependencies (\gls{war} and \gls{waw}):} To prevent data corruption, the Scoreboard must explicitly stall execution when \gls{war} or \gls{waw} hazards occur. Although out-of-order execution partially counterbalances these stalls by allowing unconstrained instructions to bypass waiting ones, name dependencies remain a major source of hardware stalls.
    \item \textbf{Bus Contention and Hardware Escalation:} Merely adding more Functional Units (\glspl{fu}) yields diminishing performance returns unless the number of internal register and result buses is increased proportionally. However, expanding the bus interconnect architecture introduces substantial hardware complexity and area overhead.
    \item \textbf{Small Issue Window and Basic Block Bounds:} The degree of exploitable \gls{ilp} is strictly bounded by the small size of the Scoreboard table (the \textit{issue window}). Moreover, basic Scoreboarding limits its scheduling logic to single basic blocks (straight-line instruction sequences without branches), preventing the processor from dynamically exploiting parallelism across branch boundaries.
\end{itemize}

\section{Implementation: Tomasulo's Algorithm}
\label{sec:tomasulo_algorithm}

As established in the previous section, basic Scoreboarding introduces critical microarchitectural bottlenecks. Operands are read strictly from the register file only when both sources become available, preventing real-time data forwarding. Furthermore, non-overlapping stage latencies between Write Result and Read Operands introduce forced stall cycles, while artificial name dependencies (\gls{war} and \gls{waw} hazards) cause severe hardware stalls that centralized issue logic cannot bypass. Merely adding extra \glspl{fu} fails to resolve these structural limits unless internal bus interconnects and issue window capacities are expanded proportionally.

To overcome these structural limitations, Robert Tomasulo designed a pioneering dynamic scheduling algorithm for the IBM System/360 Model 91 floating-point unit in 1967. Tomasulo's algorithm tracks data dependencies dynamically to minimize \gls{raw} hazards while completely eliminating \gls{war} and \gls{waw} hazards through hardware-based register renaming. By decoupling execution from architectural register names, the algorithm extracts higher levels of \gls{ilp} without requiring compiler-driven code restructuring a technique that remains a fundamental pillar of modern out-of-order processors.

\begin{definition}{Tomasulo's Algorithm}{tomasulo-algorithm}
Tomasulo's algorithm is a dynamic hardware-scheduling scheme that combines distributed dependency tracking with implicit hardware register renaming to eliminate name dependencies (\gls{war} and \gls{waw}) and bypass \gls{raw} stalls via result broadcasting over a common data bus.
\end{definition}

While built upon the fundamental dynamic scheduling principles of Scoreboarding, Tomasulo's scheme introduces four critical architectural innovations:

\begin{itemize}
\item \textbf{Distributed Control via Reservation Stations (\glspl{rs}):} Rather than relying on a centralized hardware table, control logic and operand buffering are distributed directly to the individual \glspl{fu} via sets of \glspl{rs}.
\item \textbf{Implicit Register Renaming:} Pending destination registers and instruction outputs are designated by unique reservation station tags rather than architectural register numbers. This renames registers dynamically, eliminating name dependencies (\gls{war} and \gls{waw}) at the issue phase.
\item \textbf{Unified Memory Buffering:} Load and store buffers are treated as distinct functional units with dedicated scheduling logic, allowing memory access dependencies and address calculations to be dynamically resolved.
\item \textbf{Broadcasting via the Common Data Bus (\gls{cdb}):} Results computed by any \gls{fu} are broadcast directly across a dedicated interconnect the \gls{cdb} to all waiting \glspl{rs}, load/store buffers, and register entries simultaneously, implementing automatic data forwarding and bypassing the register file writeback delay.
\end{itemize}

\subsection{\glsfmtshort{fp} Architecture Based on Tomasulo's Algorithm}
\label{subsec:tomasulo_fp_arch}

In an architecture using Tomasulo's algorithm, centralized register files are augmented by distributed hardware buffers attached directly to each \gls{fu}, known as **Reservation Stations** (\glspl{rs}). 

Each \gls{rs} entry consists of the following tracking fields:
\begin{itemize}
    \item \textbf{Tag / Name:} The unique identifier assigned to the reservation station.
    \item \textbf{Busy:} A boolean flag indicating whether the reservation station is currently occupied.
    \item \textbf{Op:} The specific operation to be executed by the attached \gls{fu} (e.g., \texttt{ADD.D}, \texttt{MULT.D}).
    \item \textbf{$V_j, V_k$:} The actual data values of the source operands. These fields are populated if the operands are ready at Issue or when they are subsequently broadcast over the \gls{cdb}.
    \item \textbf{$Q_j, Q_k$:} The tags of the reservation stations that will produce the source operands $j$ and $k$ (if the values are not yet available). If an operand is already present, its corresponding $Q$ field is cleared ($0$).
    \item \textbf{$A$:} Holds immediate values or offset addresses for memory operations (Load/Store buffers).
\end{itemize}

Additionally, the **Register Status Table** tracks pending writes to the architectural register file. Each register $R_i$ contains a field, $Q_i$, which holds the tag of the \gls{rs} that will compute its future value. If no active instruction is slated to write to $R_i$, $Q_i$ remains blank or zero.

\subsubsection{1. Issue}
Instructions are fetched sequentially from the instruction queue and issued in program order. During this stage:
\begin{enumerate}
    \item \textbf{Structural Hazard Check:} The hardware determines whether a matching \gls{rs} (or Load/Store buffer) is available (\texttt{Busy} = \textit{false}). If all matching stations are occupied, a structural hazard occurs, and the instruction stalls at Issue, blocking subsequent instructions from issuing.
    \item \textbf{Operand Fetch and Register Renaming:} If a station is free, the instruction is dispatched. The hardware checks the source registers in the Register Status Table:
    \begin{itemize}
        \item If a source register is ready ($Q_i$ is blank), its value is read directly from the register file into $V_j$ or $V_k$.
        \item If a source register is pending ($Q_i$ contains an \gls{rs} tag), that producing tag is copied into $Q_j$ or $Q_k$, dynamically renaming the operand to point to the producing station rather than the register file.
    \end{itemize}
    \item \textbf{Destination Register Update:} The destination register's entry in the Register Status Table ($Q_i$) is set to the tag of the newly allocated \gls{rs}.
\end{enumerate}

\subsubsection{2. Execute}
Once issued, an instruction waits locally at its assigned \gls{rs} for data dependencies to resolve, preventing stalls in the issue logic:
\begin{enumerate}
    \item \textbf{\gls{cdb} Monitoring:} If one or both source operands are pending ($Q_j \neq 0$ or $Q_k \neq 0$), the \gls{rs} snoops the \gls{cdb}. When a matching tag is broadcast, the \gls{rs} copies the value directly into $V_j$ or $V_k$ and clears the corresponding $Q$ field.
    \item \textbf{Execution Start:} When both source operands are available ($Q_j = 0$ and $Q_k = 0$), the operation proceeds to execution. The attached \gls{fu} computes the result over its designated latency, avoiding \gls{raw} hazard stalls.
    \item \textbf{Memory Operations:} Load and Store operations compute their effective address during execution (storing the base register + offset into field $A$) before accessing memory.
\end{enumerate}

\subsubsection{3. Write Result}
When the \gls{fu} completes computation, the result is broadcast to clear dependent instructions:
\begin{enumerate}
    \item \textbf{Broadcasting:} The \gls{fu} broadcasts the computed value and its assigned \gls{rs} tag across the \gls{cdb}.
    \item \textbf{Snooping and Updating:} All waiting \glspl{rs}, Load/Store buffers, and Register Status entries snoop the bus simultaneously:
    \begin{itemize}
        \item Any \gls{rs} with a matching $Q_j$ or $Q_k$ field captures the broadcast value into $V_j$ or $V_k$ and clears the $Q$ field.
        \item Any register status entry ($Q_i$) matching the broadcasting \gls{rs} tag writes the value to the architectural register file and clears $Q_i$.
    \end{itemize}
    \item \textbf{Deallocation:} The producing \gls{rs} resets its \texttt{Busy} flag to \textit{false}, freeing the station for subsequent instructions.
\end{enumerate}


\subsection{Architectural Trade-offs and Limitations}
\label{subsec:tomasulo_tradeoffs}

While Tomasulo's algorithm significantly outperforms Scoreboarding by enabling dynamic execution and bypassing register bottlenecks, its microarchitectural behavior introduces specific trade-offs, advantages, and physical execution limits.

\subsubsection{Dynamic Loop Unrolling and Memory Disambiguation}
A major strength of Tomasulo's algorithm is its ability to perform **implicit dynamic loop unrolling** at runtime. Because destination registers are immediately renamed to \gls{rs} tags upon issue, instructions from consecutive loop iterations can be issued into reservation stations without stalling on \gls{war} or \gls{waw} register dependencies. This allows functional units to execute operations from multiple iterations simultaneously.

However, register renaming alone does not eliminate hazards occurring through memory:
\begin{itemize}
    \item \textbf{Memory Aliasing:} Load and Store instructions may target identical memory locations out of program order, introducing potential memory-level \gls{raw}, \gls{war}, or \gls{waw} hazards.
    \item \textbf{Address Disambiguation:} To preserve execution correctness, effective address calculations are computed in order. Load and Store buffers compare target addresses; a Load operation is withheld from accessing memory until all prior Stores with matching or unresolved addresses have cleared, preventing memory aliasing faults.
\end{itemize}

\subsubsection{Summary of Architectural Advantages}
\begin{itemize}
    \item \textbf{Elimination of Name Dependencies:} Register renaming completely neutralizes \gls{war} and \gls{waw} hazards at the Issue stage, preventing artificial stalls.
    \item \textbf{Distributed Data Forwarding:} The \gls{cdb} broadcasts values directly to waiting reservation stations and load/store buffers, eliminating the latency penalty of writing back to and reading from the physical register file.
    \item \textbf{Overlapped Loop Execution:} By unrolling loops in hardware, the processor extracts maximum \gls{ilp} across iteration boundaries without requiring compiler intervention or code expansion.
\end{itemize}

\subsubsection{Key Drawbacks and Microarchitectural Bottlenecks}
\begin{itemize}
    \item \textbf{Hardware and Wiring Complexity:} Each \gls{rs} entry requires complex associative comparison hardware (snooping logic) to check broadcasted \gls{cdb} tags against pending $Q_j$ and $Q_k$ fields every clock cycle.
    \item \textbf{The \glsfmtshort{cdb} Bottleneck:} A single \gls{cdb} restricts total system throughput by allowing only one \gls{fu} to write its result per cycle. While adding multiple buses increases bandwidth, it requires multiplying input ports and comparator logic across every \gls{rs} entry and register file slot.
    \item \textbf{Control Dependency Limits:} Without speculative execution and branch prediction, instruction issue stalls at control flow boundaries (e.g., conditional branches), capping the exploitable \gls{ilp} to the size of individual basic blocks.
    \item \textbf{Imprecise Exceptions:} Because instructions complete out of program order, an exception raised by a long-latency instruction (e.g., floating-point divide) may occur after subsequent instructions have already modified the register file or memory state. This lack of in-order retirement creates imprecise exceptions, making state recovery and debugging significantly more difficult without additional hardware like a Reorder Buffer (\gls{rob}).
\end{itemize}